\appendix
\titleformat{\chapter}[display]{\centering\bfseries\fontsize{14pt}{16pt}\selectfont}{Phụ lục \thechapter}{0.3cm}{\MakeUppercase}
\renewcommand{\cftchappresnum}{PHỤ LỤC~}
\sloppy
\chapter{Cấu hình và tái lập}
Cấu hình gốc nằm trong \texttt{configs/experiments.yaml}, \texttt{configs/experiments\_v2.yaml}, \texttt{configs/experiments\_v3.yaml} và \texttt{configs/e6\_manifest\_v1.json}. Dữ liệu chính thức, model pickle và hình sinh ra không nằm đầy đủ trong Git; kiểm tra README trước khi chạy. Lệnh \texttt{python -m pytest -q} kiểm tra tích hợp; \texttt{python -m scripts.check\_v4\_audit} xác minh dữ liệu và các cặp ứng viên đã lưu. Không chạy lại test trong quá trình tái lập tài liệu.


Để tạo lại đường học trong bản sao làm việc có đủ dữ liệu và cache V2, chạy \texttt{python -m scripts.duong\_hoc}. Script dùng năm fold V3, các tập con lồng nhau với 2, 4, 6 và 8 ảnh/lớp, fit riêng PCA cho từng tập con và không dự đoán validation chính thức hoặc test. Kết quả, ID mẫu, dự đoán và checksum được lưu trong \texttt{results/duong\_hoc/}; hai hình được lưu trong \texttt{reports/latex/figures/}. Lệnh có thể ghi đè các tệp đường học đã sinh, vì vậy cần sao lưu kết quả cũ trước khi chạy lại.
